\cotefile{135}
\scanpage{17}
% transcription
pour \unsure{deviner|décrire|choisir} $P$ (Pic. d'invariants $L,N$) telle
que $P\mid\mathbb{R}$ soit splittée -- ou ce qui revient au
$=$, telle que $P$ provienne d'une cat. de Pic.
d'inv. $M,N$, \deleted{\unsure{D'où peut on en}} \inserted{On vérifie immédiatement}
que $w(\sigma)$ ne dépend pas du choix particulier
de $L\to M\mathbin{-}0$, et que sa formation est
\inserted{fonctorielle en $M,N$ et est} compatible avec \deleted{\unsure{celle des}} images inverses
par morphismes de topos. Or se plaçant
dans la situation universelle de $M\to{}_{2}N$,
\deleted{\unsure{on trouve}} soit $\mathcal{M}\mathrel{\overset{\sigma}{-}}{}_{2}\mathcal{N}$ \inserted{sur $B$}, on trouve
aisément $\operatorname{Ext}^{i}(B;\mathcal{M},\mathcal{N})=0$ pour $i>0$,
en particulier $w(\sigma)\in\operatorname{Ext}^{3}(B;\mathcal{M},\mathcal{N})\deleted{=0}$
est nul, donc $w(\sigma)$ est nul dans tous les
cas !
% archive sheet number: 16
